\def\Ponderation{35}
\def\SigleNumero{STT-1000}
\def\NomCours{Probabilités et statistique}
\def\DateEvaluation{27 avril 2023}
\def\HeureEvaluation{18h30 à 21h20}
\def\Session{Automne 2022}
\def\NomEvaluation{Examen 2}

\def\OptionsExam{11pt,answers}

% Ne rien mettre entre les accolades si l'enseignant (ou l'enseignante) est aussi le coordonnateur (ou la coordonnatrice)
\def\Coordonnatrice{}
\def\Coordonnateur{}

% Ne rien mettre entre les accolades s'il y a plus d'une section
\def\Enseignant{Julien Miron}
\def\Enseignante{}

% Ne rien mettre entre les accolades s'il s'agit d'un cours à section unique
% Mettre le nom de chacun des enseignants et enseignantes entre les accolades
% Une case à cocher sera générée pour chacune des sections
\def\SectionA{}
\def\SectionB{}
\def\SectionC{}
\def\SectionD{}

% Concerne seulement les travaux pratiques : écrire entre les accolades le nombre maximal de membres pour une équipe
\def\NombreLignes{}		

% Si vous souhaitez qu'apparaisse une déclaration d'honneur au bas de la première page, écrivez quelque chose entre les accolades
\def\EnLigne{}

\def\Directives{
\begin{enumerate}
\item Matériel autorisé:
\begin{itemize}
\item Calculatrice scientifique {\bf autorisée par la faculté des sciences et de génie};
\item Les feuillets de formules fournis avec l'examen. Ceux-ci incluent les tables de loi nécessaires. Veuillez \textbf{ne rien écrire} sur les feuillets de formules.
\end{itemize}
\item Assurez-vous que cette évaluation comporte \numquestions ~questions réparties sur \numpages{}~pages, pour un total de \numpoints{}~points.\
\item \textbf{Veuillez ne pas dégrafer votre examen.}
\item Ajoutez vos initiales dans le bas de chaque page de l'examen.
\item Veuillez déposer votre carte d'identité admise avec photo sur votre bureau.\
\item Vous devez justifier votre réponse par une démarche claire, sauf pour les questions à choix de réponse.\
\item Le verso des feuilles peut être utilisé comme brouillon, \textbf{mais ne sera pas corrigé}. Attention à ne rien dégrafer.
\end{enumerate}
}

% Veuillez indiquer le type d'évaluation entre crochets
% Examen : option à choisir pour un examen présentiel qui sera corrigé de manière classique. Un tableau des points sera affiché sur la page couverture.
% Gradescope : option à choisir pour un examen qui sera corrigé via Gradescope. Aucun tableau de pointage n'apparaîtra et il y a quelques différences au niveau de la mise en page
% TPS :
% TPF :
% Devoir :
\def\Gradescope{.} % Mode Gradescope activé (recto seulement). Mettre {} pour recto-verso.
\documentclass{Evaluation}
\begin{document}

\mbox{}
%\newpage
%%%%%%%%%%%%%%%% QUESTION %%%%%%%%%%%%%%%%

\begin{questions}
\question\label{QCMicTests}
Un intervalle de confiance de niveau 95\% pour la moyenne théorique $\mu$ est $[ 5,11 ; 9,11]$. Pour chaque énoncé qui suit, indiquez si l'énoncé est vrai, faux ou si nous n'avons pas assez d'information pour le savoir.

\medskip

\begin{parts}
    \part[1] Au seuil 5\%, on rejette $H_0: \mu=5$ pour $H_1:\mu \neq 5$.

\medskip
    \begin{oneparcheckboxes}
        \CorrectChoice Vrai 
        \choice Faux
        \choice Nous n'avons pas assez d'information pour le savoir.
    \end{oneparcheckboxes}
\medskip
 
   \part[1] Au seuil 1\%, on rejette $H_0: \mu=5$ pour $H_1:\mu \neq 5$.

 \medskip  \begin{oneparcheckboxes}
        \choice Vrai 
        \choice Faux
        \CorrectChoice Nous n'avons pas assez d'information pour le savoir.
\medskip    \end{oneparcheckboxes}
    
   \part[1] Au seuil 10\%, on rejette $H_0: \mu=5$ pour $H_1:\mu \neq 5$.

\medskip   \begin{oneparcheckboxes}
        \CorrectChoice Vrai 
        \choice Faux
        \choice Nous n'avons pas assez d'information pour le savoir.
\medskip    \end{oneparcheckboxes}
    
   \part[1] Au seuil 5\%, on rejette $H_0: \mu=8$ pour $H_1:\mu \neq 8$.

\medskip   \begin{oneparcheckboxes}
        \choice Vrai 
        \CorrectChoice Faux
        \choice Nous n'avons pas assez d'information pour le savoir.
\medskip    \end{oneparcheckboxes}
    
%   \part[1] Au seuil 1\%, on rejette $H_0: \mu=8$ pour $H_1:\mu \neq 8$.

%\medskip   \begin{oneparcheckboxes}
%        \choice Vrai 
%        \CorrectChoice Faux
%        \choice Nous n'avons pas assez d'information pour le savoir.
%    \end{oneparcheckboxes}\medskip

\end{parts}

\vspace{0.5cm}
\question [4] Nous nous intéressons aux prix des loyers des étudiants de l'Université Laval. 
Avec un échantillon aléatoire de $n=241$ étudiants, nous obtenons l'intervalle de confiance à 95\%  pour la moyenne $\mu$ suivant: $[550;1630]$. 
Cochez tous les énoncés qui sont vrais parmi les suivants. \textit{\textbf{Une réponse cochée alors qu'elle ne devrait pas l'être entraîne une pénalité de 1 point. Le minimum possible pour cette question est 0 point.}}
 
\begin{checkboxes}
\choice Il y a 95\% des étudiants qui ont un loyer entre 550 \$ et 1630 \$
\choice Si on pigeait un échantillon de 100 étudiants pour construire un nouvel intervalle de confiance à 95\% sur la moyenne, celui-ci aurait une plus faible probabilité de contenir la vraie valeur de $\mu$ que celui obtenu à partir de 241 étudiants.
\choice Si on pigeait un autre échantillon de 241 étudiants pour construire un intervalle de confiance à 95\% sur la moyenne, l'intervalle obtenu aurait nécessairement la même longueur que celui-ci.
\correctchoice Si on pigeait plusieurs échantillons de 121 étudiants dans cette population, et qu'on construisait avec chacun d'eux l'intervalle de confiance à 95\% pour $\mu$, alors environ 95\% de ces intervalles contiendraient la vraie valeur moyenne des loyers des étudiants de l'Université Laval. 
\end{checkboxes}
\vspace{0.5cm}
\question[4]
Parmi les hypothèses suivantes, \textbf{lesquelles} sont valides dans la construction de tests comme on les présente dans le cours? \textit{\textbf{Une réponse cochée alors qu'elle ne devrait pas l'être entraîne une pénalité de 1 points. Le minimum possible pour cette question est 0 point.}}

\begin{checkboxes}
\correctchoice $H_0 : \mu = 12$ vs $H_1 : \mu > 12$
\choice $H_0 : \mu \leq 12$ vs $H_1 : \mu > 12$
\choice $H_0 : \bar{x} = 4$ vs $H_1 : \bar{x} \ne 4$
\correctchoice $H_0: \mu = 1$ vs $H_1: \mu \neq 1$
\end{checkboxes}
\newpage
\question[4]
Sachant que $\displaystyle \frac{(n-1)S^2}{\sigma^2}$ suit une loi du khi-carré à $n-1$ degrés de liberté, donnez $\text{Var}(S^2)$. \textbf{\textit{Donnez les détails de votre démarche.}}

\begin{solution}
On peut calculer que 
\begin{align*}
Var(S^2) &= Var\left( \frac{\sigma^2}{(n-1)}\frac{(n-1)}{\sigma^2}S^2 \right)\\
    &= \frac{\sigma^4}{(n-1)^2} Var\left(\frac{(n-1)}{\sigma^2}S^2\right)  \\
    &= \frac{\sigma^4}{(n-1)^2}(2(n-1)) \\
    &= \frac{2\sigma^4}{(n-1)}
\end{align*}
\end{solution}

\vspace{10cm}

\question[3]  Aurélien suppose que le temps d'attente dans un bureau de la SAAQ suit une exponentielle de moyenne 240 minutes. 
Il veut construire un intervalle de confiance bilatéral à $95\%$ pour le vérifier.
Il note le temps d'attente de $16$ personnes choisies au hasard et obtient $\Bar{x}=250$ et $s^2=59~536$.
Comme il construit un intervalle de confiance, il choisit $z_{\alpha/2}=1,96$, à partir de sa table de la loi normale.
Son intervalle de confiance de niveau approximatif $95\%$ est donc $[130,44;369,56]$.
Étant donné que celui-ci contient la valeur $240$, il se dit que cette valeur est plausible pour $\mu$. 
Après avoir vérifié les résultats d'Aurélien, vous constatez qu'il y a une erreur dans ce qu'il a fait. 
Quelle est cette erreur?
\begin{solution}
    Comme ce ne sont pas des v.a. normales, pour utiliser le TCL, on doit avoit $n\geq30$, ce qui n'est pas le cas ici. Donc résulats non valides.
\end{solution}
\newpage
%%%%%%%%%%%%%%%% QUESTION %%%%%%%%%%%%%%%%
\question\label{rls}

Afin d'estimer le volume de bois disponible dans la forêt entre Burtigny et Saint-Oyens, nous aimerions modéliser le volume de bois d'un arbre à l'aide de sa circonférence. Nous disposons du volume ($Y$) et de la circonférence ($X$) de 31 arbres. Afin d'étudier le modèle $Y_i = \beta_0 + \beta_1 X_i + E_i$, en supposant que les postulats de régression linéaire sont respectés, nous avons résultats suivants:

\begin{align*}
	\sum_{i=1}^{31} x_i &= 410,70;  \quad \sum_{i=1}^{31} x_i^2 = 5\; 736,55 ; \quad \sum_{i=1}^{31} x_iy_i =13\; 887,86;\quad
	\sum_{i=1}^{31} y_i &= 935,30; \quad \sum_{i=1}^{31} y_i^2 = 36\; 324,99.
\end{align*}

\begin{parts}
	\part[7] Donnez l'équation de la droite de régression permettant de prédire le volume de bois d'un arbre de circonférence précise. \textbf{\textit{Donnez les détails de votre démarche.}}
	
	
	\begin{solution} $$ \widehat{\beta}_1 
		= \frac{S_{XY}}{S_{XX}}
		= \frac{\sum x_iy_i - n\overline{x}\overline{y}}{\sum x_i^2 - n\overline{x}^2} 
		= \frac{\displaystyle 13\; 887.86 -  \frac{410.7 \times  935.3}{31}}{\displaystyle 5\; 736.55 - \frac{410.7^2 }{ 31} }
		= \frac{1496.644}{295.4374}
		= 5.065858 $$
	
	$$ \widehat{\beta}_0 =  \overline{y}-\widehat{\beta}_1 \overline{x} 
		= \frac{935.3}{31} -  5.065858 \frac{410.7}{31} 
		= -36.94348 $$ 
		
		$$ \widehat{Y}_i = -36.94348 + 5.065858  x_i$$
		\end{solution}
	\vspace{18cm}


	\pagesuivante{\ref{rls}}

\newpage 

\suite{\ref{rls}}

	\part[3] Calculez la prévision ponctuelle du volume de bois d'un arbre ayant une circonférence de 11. \textbf{\textit{Donnez les détails de votre démarche.}}
	
	\vspace{4cm}
	
	\begin{solution}$$ \widehat{y}_i = \widehat{\beta}_0 + \widehat{\beta}_1 x_i
		= -36.94348 + 5.065858  \times 11 = 18.78096$$
		\end{solution}
	


	\part[2] Alex aimerait estimer le volume de bois de l'arbre devant sa maison. Cet arbre a une circonférence de 14. La prévision ponctuelle du volume de bois de cet arbre est environ 33,98. Choisissez l'intervalle de valeurs ayant 95\% de chances de contenir le volume de bois de cet arbre de circonférence 14.
	
	\begin{checkboxes}
			\choice $ \displaystyle	33,98 \pm t_{0,05;29} \sqrt{MS_E \left[ \frac{1}{n} + \frac{(14-\overline{x})^2}{S_{XX}}\right]} $ 
   \choice $ \displaystyle	33,98 \pm t_{0,025;1} \sqrt{MS_E \left[ \frac{1}{n} + \frac{(14-\overline{x})^2}{S_{XX}}\right]} $ 
   			\choice $ \displaystyle	33,98 \pm t_{0,025;29} \sqrt{MS_E \left[ \frac{1}{n} + \frac{(14-\overline{x})^2}{S_{XX}}\right]} $ \vspace{0.2cm}
			\correctchoice $ \displaystyle		33,98 \pm t_{0,025;29} \sqrt{MS_E \left[ 1+ \frac{1}{n} + \frac{(14-\overline{x})^2}{S_{XX}}\right]} $
   \choice $ \displaystyle		33,98 \pm t_{0,025;1} \sqrt{MS_E \left[ 1+ \frac{1}{n} + \frac{(14-\overline{x})^2}{S_{XX}}\right]} $
	\end{checkboxes}
\vspace{0.5cm}
		\part[1] Si Alex avait un arbre ayant une circonférence de 18, l'intervalle de valeurs calculé à la question précédente serait... 
	\begin{checkboxes}
			\choice plus court 
			\choice de même longueur  
			\correctchoice plus long   
		\choice on ne peut pas savoir la différence de longueur 
	\end{checkboxes}
\vspace{0.5cm}

\part[2] La valeur-p du test de significativité du lien linéaire entre $X$ et $Y$ est égale à $0,00014$. Cochez l'énoncé qui est vrai.

	\begin{checkboxes}
            \choice Il n'y a pas de lien linéaire significatif entre $X$ et $Y$.
            \choice Il y a un lien linéaire significatif entre $X$ et $Y$.
            \correctchoice Nous ne pouvons rien dire car nous ne connaissons pas le seuil du test.
            \choice Nous ne pouvons rien dire car nous ne connaissons pas $t_0$.
	\end{checkboxes}

\end{parts}
\newpage

%%%%%%%%%%%%%%%% QUESTION %%%%%%%%%%%%%%%%
\question Soit $X_1,\hdots, X_{100}$, des variables aléatoires indépendantes et identiquement distribuées selon une loi normale d’espérance $\mu$ et de variance $\sigma^2 = 400$. 
Nous désirons construire un intervalle de confiance de niveau $95\%$ pour $\mu$. Bien sûr, l'intervalle de confiance aura la forme d'un des intervalles vus au cours. Nous notons cet intervalle $[C_1;C_2]$.
\begin{parts}
    \part[3] Que vaut $\mathbb{E}[C_1]$? \textbf{\textit{Donnez les détails de votre démarche.}}
    \vspace{3cm}
    \part[3] Que vaut $\text{Var}[C_1]$? \textbf{\textit{Donnez les détails de votre démarche.}}
    \vspace{3cm}
    \part[2] Quelle est la loi de probabilité de $C_1$?
    \vspace{1.2cm}
    \part[8]Quelle est la probabilité que $C_1$ soit inférieure ou égale à $\mu$? \textbf{\textit{Donnez les détails de votre démarche.}}
    
\end{parts}


\newpage


\question[10]
Nous supposons deux populations distribuées selon des lois normales de moyennes $\mu_1$ et $\mu_2$ et de variances $\sigma^2_1$ et $\sigma^2_2$ inconnues qui sont supposées égales. Nous avons les échantillons

\begin{align*}
X_{11}, \ldots, X_{1n_1} \overset{i.i.d.}{\sim} \mathcal{N}(\mu_1, \sigma^2_1), & & X_{21}, \ldots, X_{2n_2} \overset{i.i.d.}{\sim} \mathcal{N}(\mu_2, \sigma^2_2),
\end{align*}
où $n_1=8$ et $n_2=8$. Nous avons construit un intervalle de confiance à $95\%$ pour $\mu_1-\mu_2$ et obtenu $[0,1; 4,39]$. Sachant que $S^2_1=4,52$, quelle est la valeur de $S^2_2$? \textit{\textbf{Donnez les détails de votre démarche.}}



\newpage 
\question[10] Soit $X_1,\hdots, X_{36}$, des variables aléatoires indépendantes et identiquement distribuées selon une loi normale d'espérance~$\mu$ et de variance $\sigma^2=1296$.
Nous désirons effectuer un \textbf{test unilatéral à droite} avec $H_0:\mu=10$ au seuil $\alpha = 0,025$.
Nous avons les résultats suivants:
\begin{align*}
    \bar{x}= 11,34 &  & s^2=754.
\end{align*}
Si nous savons qu'en réalité $\mu=\mu_1$ et que la puissance du test est $0,13567$, que vaut $\mu_1$? \textbf{\textit{Donnez les détails de votre démarche.}}

\newpage 



\question[10] Des chercheurs veulent comparer les traitements $A$ et $B$ contre une maladie.
Ils administrent le traitement $A$ à un groupe de 45 personnes malades et le traitement $B$ à un groupe de 36 autres personnes malades, puis ils comptent le nombre de guérisons dans chacun des groupes.
Les chercheurs pensent que le traitement $A$ est plus efficace que le traitement $B$ et veulent vérifier cette affirmation en faisant un test avec les hypothèses suivantes:

\begin{align*}
    H_0:&p_A=p_B,\\
    H_1:&p_A>p_B,
\end{align*}
où $p_A$ et $p_B$ représentent la probabilité de guérison dans les groupes $A$ et $B$ respectivement.
S'il y a eu 18 guérisons dans le groupe $B$, quel doit être le nombre minimal de guérisons dans le groupe $A$ pour que $H_0$ soit rejetée au seuil approximatif $15,866\%$? \textbf{\textit{Donnez les détails de votre démarche.}}

\vspace{0.5cm}

\textit{Coup de pouce:} $\displaystyle \hat{p}(1-\hat{p})=\frac{-25 \hat{p}_A^2+25 \hat{p}_A+14}{81}$.
\newpage


\question Supposons que les variables aléatoires $X_1,...,X_{20}$ suivent une loi exponentielle décalée de paramètre $\theta \geq0 $, dont la \textbf{fonction de densité} est donnée par 

 \begin{equation*}
  f_X(x) = \begin{cases}
         e^{-(x-\theta)},\ \text{  si  } \theta \leq x , \\
        0, \text{ sinon.}
        \end{cases}
 \end{equation*}

Nous aimerions faire un test statistique pour savoir si $\theta=0$. Nous avons donc les hypothèses

\begin{table}[ht]
    \centering
    \begin{tabular}{cc}
         $H_0: \theta=0$ & $H_1: \theta > 0$.
    \end{tabular}
\end{table}
La statistique du test est
\begin{align*}
    M=\min (X_1,...,X_{20}),
\end{align*}
qui correspond à la valeur minimale de notre échantillon (nous pouvons aussi écrire $M=X_{(1)}$, dans la notation introduite au cours).
Nous savons que sous $H_0$, $M$ suit une loi exponentielle de paramètre $n=20$. Sa \textbf{fonction de répartition} est donnée par 

 \begin{equation*}
  F_M(m) = \begin{cases}
         1-e^{-20m},\ \text{  si  } 0\leq m \\
        0, \text{ sinon.}
        \end{cases}
 \end{equation*}

La table suivante donne quelques valeurs de cette fonction de répartition.
\begin{table}[ht]
    \centering
\begin{tabular}{|cc|}
  \hline
 $m$ & $F_M(m)$ \\ 
  \hline
    0,0005 & 0,01 \\
    0,0025 & 0,05 \\
    0,0053 & 0,10 \\
    0,1151 & 0,90 \\
    0,1498 & 0,95 \\ 
    0,2303 & 0,99 \\ 
  \hline
\end{tabular}
\end{table}

Nous avons l'échantillon $x_1,\hdots,x_{20}$ suivant:

\begin{table}[ht]
\centering
\begin{tabular}{rrrrrrrrrr}
  \hline
0,1613 & 0,1933 & 0,2980 & 0,3331 & 0,3423 & 0,4392 & 0,4472 & 0,4580 & 0,5289 & 0,5530 \\ 
0.5570 & 0.7737 & 0.8574 & 0.9376 & 1.2509 & 1.4475 & 1.9064 & 2.6019 & 2.8685 & 2.9829 \\ 
   \hline
\end{tabular}
\end{table}
\begin{parts}
\part[3] Au seuil $\alpha=5\%$, quelle est la \textit{\textbf{décision}} du test. \textbf{\textit{Une décision sans justification n'obtiendra aucun point, même si la décision est bonne.}}
\vspace{3cm}
\part[3] \textbf{Calculez} la valeur-p du test en ne gardant que 4 chiffres après la virgule. \textbf{\textit{Donnez les détails de votre démarche.}}
\end{parts}


\end{questions}

\end{document} 
